% © 2025 Ing. David Jaroš — CC BY-NC-ND 4.0
%
% This work is licensed under a Creative Commons Attribution-NonCommercial-NoDerivatives
% 4.0 International License (CC BY-NC-ND 4.0).
%
% License History: Earlier drafts (up to v0.3) were released under CC BY 4.0.
% From v0.4 onward, all material is released under CC BY-NC-ND 4.0 to protect
% the integrity of the theoretical work during ongoing academic development.
%
% See LICENSE.md for full license text.

\documentclass[12pt,a4paper]{article}
\usepackage{amsmath, amssymb, amsthm}
\usepackage{geometry}
\usepackage{hyperref}
\usepackage{booktabs}
\usepackage{tcolorbox}
\geometry{a4paper, margin=1in}

\newtheorem{theorem}{Theorem}[section]
\newtheorem{proposition}[theorem]{Proposition}
\newtheorem{lemma}[theorem]{Lemma}
\newtheorem{corollary}[theorem]{Corollary}
\newtheorem{definition}[theorem]{Definition}
\newtheorem{axiom}{Axiom}
\newtheorem{remark}{Remark}[section]

% Proof-level macros
\newcommand{\Lzero}{\textbf{[L0]}}
\newcommand{\Lone}{\textbf{[L1]}}
\newcommand{\Ltwo}{\textbf{[L2]}}
\newcommand{\SE}{\textbf{[SE]}}
\newcommand{\PROVED}{\textbf{PROVED}}

\title{\textbf{Unified Biquaternion Theory}\\[6pt]
\large A Referee-Ready Presentation}
\author{Ing.\ David Jaroš}
\date{2025}

\begin{document}
\maketitle

\begin{abstract}
The Unified Biquaternion Theory (UBT) presents a framework for unifying
General Relativity, Quantum Electrodynamics (QED), and Standard Model gauge
symmetries within a single biquaternionic field $\Theta(q,\tau)$ defined over
complex time $\tau = t + i\psi$.  The theory is built on three structural
axioms and an action principle.  Key results include: (i) exact recovery of
Einstein's field equations in the real-projection limit; (ii) derivation of
the QED Lagrangian and Maxwell's equations from U(1) gauge invariance of
$\Theta$; (iii) emergence of the gauge group
$\mathrm{SU}(3)_c\times\mathrm{SU}(2)_L\times\mathrm{U}(1)_Y$
from the algebraic structure $\mathbb{H}\otimes_\mathbb{R}\mathbb{C}$;
(iv) three-generation structure from imaginary-time winding modes.
Open problems — including the fine-structure-constant exponent, the Weinberg
angle, and fermion mass parameters — are stated explicitly and without
overstatement.  Predictions distinguishing UBT from the Standard Model are
identified and assessed for experimental falsifiability.
\end{abstract}

\noindent\textbf{Keywords}: biquaternion, complex time, general relativity,
QED, Standard Model, gauge symmetry, fermion mass, fine structure constant.

\tableofcontents
\newpage

% ======================================================================
\section{Introduction}
% ======================================================================

\subsection{Motivation}

Modern fundamental physics rests on two pillars — General Relativity (GR)
and Quantum Field Theory (QFT) — that are logically incompatible at high
energies.  Numerous proposals for their unification exist (string theory,
loop quantum gravity, non-commutative geometry, twistor theory), but none
has achieved experimental confirmation or mathematical completeness.

UBT takes a different starting point: the observation that the biquaternion
algebra $\mathcal{B} = \mathbb{H}\otimes_\mathbb{R}\mathbb{C}$, the simplest
associative algebra combining real quaternions with complex numbers, contains
the algebraic seeds of both GR geometry and gauge symmetries.
By placing the fundamental field $\Theta$ in $\mathcal{B}$ and extending
time to the complex plane $\tau = t + i\psi$, GR and the Standard Model
gauge groups emerge as projections of a single mathematical object.

\subsection{Scope of This Paper}

This paper is organized as a referee-ready presentation of the core UBT
results with explicit proof-level labels:
\begin{itemize}
  \item \Lzero{} — algebraic identity, exact, no free parameters
  \item \Lone{} — proved given the stated axioms
  \item \Ltwo{} — requires additional lemmas not yet proved
  \item \SE{} — semi-empirical: formula present but parameters fitted
  \item $[\mathrm{O}]$ — open problem
\end{itemize}

\noindent Detailed derivations are in the companion canonical documents
(see \texttt{canonical/bridges/} and \texttt{canonical/interactions/}).

\subsection{Plan of the Paper}

\begin{itemize}
  \item \S\ref{sec:axioms}: Axioms and foundational assumptions
  \item \S\ref{sec:algebra}: Biquaternion algebra
  \item \S\ref{sec:field}: Theta field definition
  \item \S\ref{sec:action}: Action principle
  \item \S\ref{sec:GR}: General Relativity recovery
  \item \S\ref{sec:QED}: QED limit and Maxwell equations
  \item \S\ref{sec:cosmology}: Cosmology and dark sector
  \item \S\ref{sec:predictions}: Testable predictions
  \item \S\ref{sec:open}: Open problems
\end{itemize}

% ======================================================================
\section{Axioms}
\label{sec:axioms}
% ======================================================================

UBT is built on a minimal set of structural axioms.  They are not derived
but define the theory; all results below follow from them.

\begin{axiom}[Biquaternion Algebra]\label{ax:algebra}
The fundamental mathematical arena is:
\begin{equation}
  \mathcal{B} \;:=\; \mathbb{H}\otimes_{\mathbb{R}}\mathbb{C}
  \;\cong\; \mathrm{Mat}(2,\mathbb{C}).
\end{equation}
\textit{Source}: \texttt{canonical/fields/biquaternion\_algebra.tex}, \S1.
\end{axiom}

\begin{axiom}[Fundamental Field]\label{ax:field}
The fundamental dynamical object is:
\begin{equation}
  \Theta \;:\; \mathbb{R}^{1,3}\times\mathbb{C} \;\to\; \mathcal{B},
  \qquad (q,\tau) \;\mapsto\; \Theta(q,\tau).
\end{equation}
All physical fields (metric, gauge fields, matter) are encoded in $\Theta$.
\textit{Source}: \texttt{canonical/fields/theta\_field.tex}, \S1.
\end{axiom}

\begin{axiom}[Complex Time]\label{ax:time}
Time is extended to the complex plane:
\begin{equation}
  \tau \;=\; t + i\psi, \qquad t,\psi\in\mathbb{R}.
\end{equation}
Real observables correspond to the section $\psi = 0$.
\textit{Source}: \texttt{canonical/fields/biquaternion\_time.tex}, \S1.
\end{axiom}

\begin{axiom}[Field Equation]\label{ax:eq}
The field $\Theta$ satisfies the biquaternionic wave equation:
\begin{equation}\label{eq:tshirt}
  \boxed{\nabla^\dagger\nabla\,\Theta(q,\tau) \;=\; \kappa\,\mathcal{T}(q,\tau),}
\end{equation}
where $\nabla^\dagger\nabla$ is the biquaternionic Laplacian,
$\mathcal{T}$ is the source tensor, and $\kappa = 8\pi G/c^4$.
\textit{Source}: \texttt{canonical/fields/theta\_field.tex}, \S3.
\end{axiom}

\begin{axiom}[Real Projection]\label{ax:proj}
Physical observables are recovered by the real projection:
\begin{equation}
  g_{\mu\nu} \;=\; \mathrm{Re}(\mathcal{G}_{\mu\nu}[\Theta]),
\end{equation}
where $\mathcal{G}_{\mu\nu}$ is the biquaternionic metric.
\textit{Source}: \texttt{canonical/geometry/metric.tex}, \S2.
\end{axiom}

\begin{axiom}[Action Principle]\label{ax:action}
The dynamics are governed by $\delta S = 0$ with:
\begin{equation}\label{eq:S}
  S[\Theta] \;=\; \int d^4x\,\sqrt{-g}\,
  \Bigl[
    \tfrac{1}{2\kappa}R[\Theta]
    + \mathcal{L}_{\mathrm{kin}}[\Theta]
    - \tfrac{1}{4}F_{\mu\nu}^a F^{a\mu\nu}
    - V(\Theta)
  \Bigr].
\end{equation}
\textit{Source}: \texttt{THEORY/canonical/canonical\_action.tex}.
\end{axiom}

% ======================================================================
\section{Biquaternion Algebra}
\label{sec:algebra}
% ======================================================================

\subsection{Structure}

The biquaternion algebra $\mathcal{B} = \mathbb{H}\otimes_\mathbb{R}\mathbb{C}$
has basis $\{1, i, j, k\} \times \{1, \mathbf{i}\}$, where $i,j,k$ are the
quaternion units and $\mathbf{i}$ is the complex imaginary.
Key relations:
\begin{equation}
  i^2 = j^2 = k^2 = ijk = -1, \qquad \mathbf{i}^2 = -1.
\end{equation}
Every element $Q\in\mathcal{B}$ can be written as $Q = q_0 + q_1 i + q_2 j + q_3 k$
with $q_\mu \in \mathbb{C}$.

\begin{proposition}[Matrix representation, \Lzero]
$\mathcal{B} \cong \mathrm{Mat}(2,\mathbb{C})$ via:
\begin{equation}
  i \mapsto \begin{pmatrix} \mathbf{i} & 0 \\ 0 & -\mathbf{i} \end{pmatrix},\quad
  j \mapsto \begin{pmatrix} 0 & 1 \\ -1 & 0 \end{pmatrix},\quad
  k \mapsto \begin{pmatrix} 0 & \mathbf{i} \\ \mathbf{i} & 0 \end{pmatrix}.
\end{equation}
\end{proposition}

This isomorphism is exact.  It implies that $\mathcal{B}$ contains all
$2\times 2$ complex matrices, which is the representation space for the
Standard Model matter fields.

\subsection{Involution Structure}

$\mathcal{B}$ has a $\mathbb{Z}_2\times\mathbb{Z}_2\times\mathbb{Z}_2$
involution structure arising from quaternionic conjugation, complex conjugation,
and their composition.  The three involutions select the subgroups
$\mathrm{SU}(3)_c$, $\mathrm{SU}(2)_L$, and $\mathrm{U}(1)_Y$ from the
automorphism group of $\mathcal{B}$.

\textit{Source}: \texttt{canonical/algebra/involutions\_Z2xZ2xZ2.tex};
\texttt{canonical/bridges/gauge\_emergence\_bridge.tex}.

% ======================================================================
\section{The Theta Field}
\label{sec:field}
% ======================================================================

\subsection{Definition}

\begin{definition}[Theta field]
The fundamental field is a map
$\Theta : \mathbb{R}^{1,3}\times\mathbb{C} \to \mathrm{Mat}(2,\mathbb{C})$.
In explicit form, with complex time $\tau = t+i\psi$:
\begin{equation}
  \Theta(q,\tau) \;=\;
  \begin{pmatrix}
    \Theta_{11}(q,\tau) & \Theta_{12}(q,\tau) \\
    \Theta_{21}(q,\tau) & \Theta_{22}(q,\tau)
  \end{pmatrix}, \qquad
  \Theta_{ij} \in \mathbb{C}.
\end{equation}
\end{definition}

\subsection{Field Content}

The matrix structure encodes:
\begin{itemize}
  \item \textbf{Real diagonal}: gravitational (metric) sector
  \item \textbf{Imaginary diagonal}: U(1) phase (electromagnetic sector)
  \item \textbf{Off-diagonal}: SU(2) spinor degrees of freedom
  \item \textbf{Complex-time dependence}: SU(3) colour and generation structure
\end{itemize}

\subsection{Biquaternionic Laplacian}

The operator $\nabla^\dagger\nabla$ in the master equation \eqref{eq:tshirt} is:
\begin{equation}
  \nabla^\dagger\nabla \;=\;
  -\partial_\tau^*\partial_\tau - g^{\mu\nu}D_\mu D_\nu,
\end{equation}
where $D_\mu = \partial_\mu + \Gamma_\mu + igA_\mu$ includes the gravitational
spin connection $\Gamma_\mu$ and the gauge connection $A_\mu$.
In the limit $\psi \to 0$ this reduces to the standard covariant d'Alembertian.

\textit{Source}: \texttt{canonical/explanation\_of\_nabla.tex};
\texttt{canonical/fields/theta\_field.tex}, \S2.

% ======================================================================
\section{Action Principle}
\label{sec:action}
% ======================================================================

\subsection{Full Action}

The complete UBT action (Axiom~\ref{ax:action}) decomposes as:
\begin{equation}
  S[\Theta] = S_{\mathrm{grav}} + S_{\mathrm{kin}} + S_{\mathrm{gauge}} + S_{\mathrm{pot}},
\end{equation}
with:
\begin{align}
  S_{\mathrm{grav}} &= \int d^4x\,\frac{\sqrt{-g}}{2\kappa}\,R, \\
  S_{\mathrm{kin}} &= \int d^4x\,\sqrt{-g}\,
    \mathrm{Tr}\bigl[(D_\mu\Theta)^\dagger(D^\mu\Theta)\bigr], \\
  S_{\mathrm{gauge}} &= -\int d^4x\,\frac{\sqrt{-g}}{4}\,F_{\mu\nu}^a F^{a\mu\nu}, \\
  S_{\mathrm{pot}} &= -\int d^4x\,\sqrt{-g}\,V(\Theta).
\end{align}

\subsection{Variational Equations}

Varying with respect to $g^{\mu\nu}$ gives the Einstein equations (see \S\ref{sec:GR}).
Varying with respect to $A_\nu^a$ gives Yang--Mills equations; in the U(1) sector
these reduce to Maxwell's equations (see \S\ref{sec:QED}).
Varying with respect to $\Theta^\dagger$ gives the matter equation
$D_\mu D^\mu\Theta = \partial V/\partial\Theta^\dagger$,
which reduces to the Dirac equation for spinor modes.

\textit{Source}: \texttt{THEORY/canonical/canonical\_action.tex}, \S2.

% ======================================================================
\section{General Relativity Recovery}
\label{sec:GR}
% ======================================================================

\subsection{Chain Summary}

\begin{tcolorbox}[colback=green!5!white,colframe=green!50!black,title=GR Recovery Status]
Steps 1--5: \textbf{PROVED [L1]} (on-shell, real projection)\\
Step 6 (off-shell Θ-only closure): \textbf{OPEN [L2]}\\
\textit{The on-shell recovery fully validates GR compatibility.}
\end{tcolorbox}

The five-step chain from $\Theta$ to Einstein's equations is:

\begin{enumerate}
  \item $\Theta(q,\tau)$ defines the biquaternionic geometry (\Lzero)
  \item $g_{\mu\nu} = \mathrm{Re}(\mathcal{G}_{\mu\nu}[\Theta])$
        is the physical metric (\Lone)
  \item $\Gamma^\lambda_{\mu\nu} = \mathrm{Re}(\Omega^\lambda_{\mu\nu})$
        is the Levi-Civita connection (\Lone)
  \item $R_{\mu\nu} = \mathrm{Re}(\mathcal{R}_{\mu\nu})$
        is the Ricci tensor (\Lone)
  \item $G_{\mu\nu} = \kappa T_{\mu\nu}$ follows from variation of $S_{\mathrm{grav}}$
        (\Lone)
\end{enumerate}

\subsection{Theorem: GR is a Real Projection of UBT}

\begin{theorem}[GR Recovery, \Lone]
In the limit $\psi \to 0$ (real time, $\mathrm{Im}(\mathcal{G}_{\mu\nu}) = 0$),
the field equations \eqref{eq:tshirt} reduce exactly to Einstein's equations:
\begin{equation}
  G_{\mu\nu} \;=\; R_{\mu\nu} - \tfrac{1}{2}g_{\mu\nu}R \;=\; 8\pi G\,T_{\mu\nu}.
\end{equation}
\end{theorem}

\noindent This recovery holds for all curvature regimes (flat, weak, strong,
cosmological).  Every experimental confirmation of GR (perihelion precession,
gravitational waves, black hole imaging) therefore automatically validates the
GR sector of UBT.

\textit{Source}: \texttt{canonical/bridges/GR\_chain\_bridge.tex};
\texttt{canonical/geometry/gr\_as\_limit.tex}.

\subsection{Beyond GR: Imaginary Curvature Sector}

The imaginary part $\mathrm{Im}(\mathcal{G}_{\mu\nu})$ — the phase curvature
— is invisible to classical measurements (ordinary matter couples only to
$g_{\mu\nu} = \mathrm{Re}(\mathcal{G}_{\mu\nu})$).
It may carry dark-sector energy (see \S\ref{sec:cosmology}).

% ======================================================================
\section{QED Limit}
\label{sec:QED}
% ======================================================================

\begin{tcolorbox}[colback=green!5!white,colframe=green!50!black,title=QED Sector Status]
\textbf{PROVED [L0]}: U(1) gauge symmetry, $F_{\mu\nu}$ definition, Bianchi identity\\
\textbf{PROVED [L1]}: QED Lagrangian, Maxwell equations, Dirac equation\\
\textbf{SEMI-EMPIRICAL [SE]}: Precise value of $\alpha$ (formula structure proved;
$B_{\mathrm{base}}$ exponent open)
\end{tcolorbox}

\subsection{U(1) Gauge Symmetry}

The scalar phase right action on $\mathcal{B}$ defines an exact U(1) symmetry
(\Lzero):
\begin{equation}
  \Theta(x) \;\mapsto\; \Theta(x)\,e^{i\alpha(x)}.
\end{equation}
This requires the gauge potential $A_\mu = e^{-1}\partial_\mu\varphi$
and covariant derivative $D_\mu = \partial_\mu - ieA_\mu$.

\textit{Source}: \texttt{canonical/interactions/sm\_gauge.tex}, Theorem G.3.

\subsection{Electromagnetic Field Tensor}

\begin{equation}
  F_{\mu\nu} \;=\; \partial_\mu A_\nu - \partial_\nu A_\mu.
\end{equation}
The Bianchi identity $\partial_{[\lambda}F_{\mu\nu]} = 0$ follows algebraically
(\Lzero) and gives the homogeneous Maxwell equations.

\subsection{QED Lagrangian}

\begin{equation}
  \mathcal{L}_{\mathrm{QED}}
  \;=\;
  \mathrm{Tr}\bigl[(D_\mu\Theta)^\dagger(D^\mu\Theta)\bigr]
  \;-\;
  \frac{1}{4}\,F_{\mu\nu}F^{\mu\nu},
\end{equation}
where $D_\mu = \partial_\mu - ieA_\mu$.

\textit{Source}: \texttt{canonical/interactions/qed.tex}, eq.\ (1).

\subsection{Maxwell Equations}

Varying $\mathcal{L}_{\mathrm{QED}}$ with respect to $A_\nu$:
\begin{align}
  \partial_\mu F^{\mu\nu} &= j^\nu, \label{eq:maxwell_inhom} \\
  \partial_{[\lambda}F_{\mu\nu]} &= 0. \label{eq:maxwell_hom}
\end{align}
Together these are all four Maxwell equations (\Lone).

\textit{Source}: \texttt{canonical/bridges/Maxwell\_complete\_derivation.tex}.

\subsection{Dirac Equation}

In the U(1) sector with constant phase ($\partial_\psi = 0$), the spinor
modes of $\Theta$ satisfy:
\begin{equation}
  (i\gamma^\mu D_\mu - m)\,\psi \;=\; 0, \qquad
  D_\mu = \partial_\mu - ie A_\mu.
\end{equation}
This follows from varying $\mathcal{L}_{\mathrm{QED}}$ with respect to $\bar\psi$ (\Lone).

\textit{Source}: \texttt{canonical/bridges/QED\_complete\_chain.tex};
\texttt{unified\_biquaternion\_theory/ubt\_appendix\_10\_qed\_dirac.tex}.

\subsection{Fine Structure Constant}

The UBT formula for $\alpha$ has the structure:
\begin{equation}
  \alpha \;=\; \frac{B_0}{B_{\mathrm{base}} + B_\alpha},
\end{equation}
with $B_0 = 8\pi$ (PROVED [L1]) and $B_{\mathrm{base}} + B_\alpha \approx 87.87$
reproducing $\alpha \approx 1/137.036$.  The exponent in
$B_{\mathrm{base}} = N_{\mathrm{eff}}^{3/2}$ is currently unproved (Open Problem,
see \S\ref{sec:open}).

\textit{Source}: \texttt{research/alpha\_derivation\_status.md};
\texttt{canonical/bridges/QED\_limit\_bridge.tex}.

\subsection{Gauge Group Emergence}

\begin{tcolorbox}[colback=green!5!white,colframe=green!50!black,title=Gauge Group Status]
\textbf{PROVED [L0]}: $\mathrm{SU}(3)_c$, $\mathrm{SU}(2)_L$, $\mathrm{U}(1)_Y$
from $\mathbb{C}\otimes\mathbb{H}$ involutions\\
\textbf{SEMI-EMPIRICAL}: Weinberg angle $\sin^2\theta_W \approx 0.231$\\
\textbf{MOTIVATED (not proved)}: Chirality selection $\mathrm{SU}(2)_L$ over $\mathrm{SU}(2)_R$
\end{tcolorbox}

The three gauge groups emerge from the involution structure of $\mathcal{B}$:
\begin{itemize}
  \item $\mathrm{SU}(3)_c$ from $\mathbb{Z}_2\times\mathbb{Z}_2\times\mathbb{Z}_2$ involutions (\Lzero)
  \item $\mathrm{SU}(2)_L$ from norm-preserving left multiplications (\Lzero)
  \item $\mathrm{U}(1)_Y$ from scalar phase right multiplications (\Lzero)
\end{itemize}

\textit{Source}: \texttt{canonical/bridges/gauge\_emergence\_bridge.tex}.

% ======================================================================
\section{Cosmology}
\label{sec:cosmology}
% ======================================================================

\subsection{Dark Energy from Phase Curvature}

The imaginary part of the biquaternionic metric,
$\Im(\mathcal{G}_{\mu\nu}) = \psi_{\mu\nu}$ (phase curvature),
satisfies field equations but does not couple to ordinary (baryonic) matter.
It provides a natural candidate for the observed accelerated cosmological
expansion without fine-tuning the cosmological constant.

The effective cosmological constant in UBT arises as:
\begin{equation}
  \Lambda_{\mathrm{eff}} \;=\; \frac{8\pi G}{c^4}\,
  \rho_{\mathrm{vac}}[\psi_{\mu\nu}],
\end{equation}
where $\rho_{\mathrm{vac}}$ is the vacuum energy density of the phase-curvature
field, which is a geometric property of $\Theta$ rather than a sum of
zero-point fluctuations.

\textit{Source}: \texttt{unified\_biquaternion\_theory/priority\_P6\_dark\_energy/};
\texttt{consolidation\_project/appendix\_M\_dark\_energy\_UBT.tex}.

\subsection{Dark Matter as Topological Hopfions}

Topologically stable configurations of $\Theta$ — Hopfion solitons with
non-zero Hopf invariant $Q_H\in\mathbb{Z}$ — are candidates for dark matter.
These arise as non-trivial mappings $S^3 \to S^2$ in the field space of $\Theta$.

Hopfion dark matter candidates:
\begin{itemize}
  \item Gravitationally coupled (invisible to electromagnetic sector)
  \item Mass scale linked to internal-mode physics (correlated with $\alpha$ sector)
  \item Cold dark matter spectrum consistent with $\Lambda$CDM bounds
\end{itemize}

\textit{Source}: \texttt{consolidation\_project/appendix\_U\_dark\_matter\_unified\_padic.tex};
\texttt{unified\_biquaternion\_theory/solution\_P5\_dark\_matter/}.

\subsection{Three Generations from Winding Modes}

The three fermion generations arise from three independent $\psi$-winding modes
on the imaginary-time circle $\psi \sim \psi + 2\pi$:
\begin{itemize}
  \item Mode $n = 1, 2, 3$ → Generation 1, 2, 3
  \item Each mode carries identical gauge quantum numbers
  \item Status: \textbf{PROVED [L0]} — algebraic identity
\end{itemize}

\textit{Source}: \texttt{STATUS\_THEORY\_ASSESSMENT.md} v61;
\texttt{DERIVATION\_INDEX.md}.

\subsection{Cosmological Predictions}

UBT predicts modifications to the standard $\Lambda$CDM cosmology:
\begin{itemize}
  \item Effective number of relativistic species: $N_{\mathrm{eff}} \approx 3.09$
        (UBT prediction; current measurement $2.99\pm 0.17$)
  \item CMB acoustic peak modifications from phase curvature
  \item Gravitational wave propagation: possible frequency-dependent dispersion
\end{itemize}

Detailed predictions and their current observational status are in
\texttt{docs/predictions/UBT\_cosmology\_predictions.md}.

% ======================================================================
\section{Testable Predictions}
\label{sec:predictions}
% ======================================================================

\subsection{Predictions Summary}

The following table lists the key UBT predictions with their proof level
and current experimental status.

\begin{center}
\begin{tabular}{p{5cm}llp{3.5cm}}
\toprule
\textbf{Prediction} & \textbf{Level} & \textbf{Value} & \textbf{Exp.\ Status} \\
\midrule
GR recovery (all regimes) & \Lone & Exact & All GR tests passed \\
$\mathrm{SU}(3)_c\times\mathrm{SU}(2)_L\times\mathrm{U}(1)_Y$ & \Lzero & Algebraic & Confirmed by SM \\
Three SM fermion generations & \Lzero & 3 & Confirmed \\
Massless photon & \Lone & $m_\gamma = 0$ & Confirmed \\
Electron mass $m_e$ & \SE & $0.511$ MeV & 0.22\% accuracy \\
Fine structure constant $\alpha$ & \SE & $1/137.036$ & Verified with fitted $B_\text{base}$ \\
$N_\text{eff} \approx 3.09$ & \Ltwo & $3.090$ & Consistent (Planck: $2.99\pm0.17$) \\
Hopfion dark matter & $[\mathrm{O}]$ & Mass TBD & Not tested \\
Gravitational wave dispersion & $[\mathrm{O}]$ & Signal TBD & Future LISA/ET \\
\bottomrule
\end{tabular}
\end{center}

\subsection{Key Falsification Criterion}

UBT is falsified if any of the following are observed:
\begin{enumerate}
  \item Violation of GR in the real-projection regime (since GR recovery is exact)
  \item A fourth SM gauge group factor inconsistent with
        $\mathbb{H}\otimes_\mathbb{R}\mathbb{C}$
  \item More or fewer than three fermion generations
        (since the three-generation result is proved from winding modes)
\end{enumerate}

\textit{Source}: \texttt{TESTABILITY\_AND\_FALSIFICATION.md}.

% ======================================================================
\section{Open Problems}
\label{sec:open}
% ======================================================================

The following open problems are the current research frontier.  They are
stated clearly so that a referee can assess the theory's completeness.

\begin{enumerate}

\item \textbf{Fine structure constant $B_\text{base}$ exponent} \hfill $[\mathrm{O}]$\\
The formula $B_\text{base} = N_\text{eff}^{3/2}$ has the exponent $3/2$
unproved from first principles.  22 approaches have been tested; none succeeded.
Promising direction: Kac--Moody algebras at level $k = 1$.\\
\textit{Source}: \texttt{research/alpha\_derivation\_status.md};
\texttt{canonical/bridges/QED\_limit\_bridge.tex}, \S3.

\item \textbf{Weinberg angle} $\sin^2\theta_W \approx 0.231$ \hfill \SE{}\\
Cannot be derived from $\mathbb{C}\otimes\mathbb{H}$ alone.
Additional input from fermion representations or Higgs sector is required.\\
\textit{Source}: \texttt{canonical/bridges/gauge\_emergence\_bridge.tex}, \S3.

\item \textbf{Chirality selection} ($\mathrm{SU}(2)_L$ not $\mathrm{SU}(2)_R$)
      \hfill \Ltwo{}\\
Motivated by $\psi$-parity under $\psi\to-\psi$, but not yet elevated
to a theorem.

\item \textbf{Fermion mass parameters} \hfill \SE{}\\
The mass parameter $B_m \approx -14.099$ MeV is fitted to $m_e$, not derived.
Lepton mass ratios and quark masses are not derived.\\
\textit{Source}: \texttt{research/fermion\_mass\_program.md}.

\item \textbf{Off-shell GR closure} \hfill $[\mathrm{L2}]$\\
The purely off-shell reduction
$\mathrm{Re}(\nabla^\dagger\nabla\Theta) \to G_{\mu\nu}$
has a structural obstruction that prevents simple closure.
This does not affect the on-shell result (Steps 1--5 of the GR chain).\\
\textit{Source}: \texttt{canonical/bridges/GR\_chain\_bridge.tex}, Step 6.

\item \textbf{Geometric correction factor $R \approx 1.114$} \hfill $[\mathrm{O}]$\\
The one-loop correction in the $\alpha$ formula involves a factor
$R \approx 1.114$ of unknown geometric origin.  Eight approaches tested.

\item \textbf{Neutrino masses and PMNS matrix} \hfill $[\mathrm{O}]$\\
The neutrino sector is referenced in the repository but without a complete
mass derivation.\\
\textit{Source}: \texttt{research/fermion\_mass\_program.md}, Open Problem 4.

\item \textbf{CKM matrix} \hfill \SE{}\\
Quark mixing angles are partially treated but the full CKM matrix is not
derived.\\
\textit{Source}: \texttt{consolidation\_project/appendix\_QA2\_quarks\_CKM.tex}.

\end{enumerate}

% ======================================================================
\section*{Conclusion}
% ======================================================================

The Unified Biquaternion Theory provides a mathematically consistent
framework in which General Relativity, QED, and the Standard Model gauge
symmetries emerge as projections of a single biquaternionic field.
The core results — GR recovery, QED Lagrangian, Maxwell equations,
Dirac equation, and three-generation structure — are proved at level L1 or
better.

Significant open problems remain, particularly the first-principles derivation
of the fine-structure constant (B_base exponent), the Weinberg angle, and
fermion masses.  These are stated without overstatement in the companion
documents in \texttt{research/}.

The theory is falsifiable: it makes sharp predictions about three fermion
generations, gauge group structure, and GR recovery that are not parameterically
adjustable.

\section*{Acknowledgments}

The author acknowledges A.H.\ Chamseddine (2025) and E.\ Verlinde (2025).

\bigskip
\noindent\textbf{Companion documents}:
\begin{itemize}
  \item Complete QED chain: \texttt{canonical/bridges/QED\_complete\_chain.tex}
  \item Maxwell derivation: \texttt{canonical/bridges/Maxwell\_complete\_derivation.tex}
  \item GR chain: \texttt{canonical/bridges/GR\_chain\_bridge.tex}
  \item Gauge emergence: \texttt{canonical/bridges/gauge\_emergence\_bridge.tex}
  \item $\alpha$ status: \texttt{research/alpha\_derivation\_status.md}
  \item Fermion mass program: \texttt{research/fermion\_mass\_program.md}
\end{itemize}

\section*{License}

\noindent © 2025 Ing.\ David Jaroš — CC BY-NC-ND 4.0

This work is licensed under a Creative Commons
Attribution-NonCommercial-NoDerivatives 4.0 International License.

\textbf{License History:} Earlier drafts (up to v0.3) were released under
CC BY 4.0. From v0.4 onward, all material is released under CC BY-NC-ND 4.0.

\end{document}
